% GENERATED FILE. Edit canonical lessons, then run npm run generate:books.
% canonical-source-hash: fnv1a64:e1ea81548fb85756
% canonical-lessons: RU-C29-head, RU-C29-hand, RU-C29-leg, RU-C29-belly, RU-C29-tooth

\chapter{Head, Hand, Leg, Belly, Tooth}
\label{ch:ru-head-hand-leg-belly-tooth}
\hlchaptermodality{29}

\begin{chapteropening}
\textbf{By the end of this chapter:} \emph{I can name the head, a hand, a leg, the belly and a tooth.}

Complete the last lesson of chapter 29: I can name the head, a hand, a leg, the belly and a tooth.
\end{chapteropening}

\section[\texorpdfstring{\ru{голова} (golová)}{голова (golová)}]{\ru{голова} (golová) — the head}

\label{lesson:RU-C29-head}



\begin{quote}
Before the new one: say the Russian for first, then the Russian for second.
\end{quote}

\subsection*{You'll want to know: \ru{голова}}
\textbf{\ru{голова}} (\emph{golová}) — \textquotedblleft{}the head\textquotedblright{}.

\begin{cousinweb}
\textbf{\ru{голова}} $\leftarrow$ Proto-Slavic \emph{*golva}, \textquotedblleft{}head\textquotedblright{} — inherited, and a cousin of Lithuanian \emph{galvà}, \textquotedblleft{}head\textquotedblright{} (Vasmer). East Slavic says \emph{-\ru{оло}-} where Church Slavonic says \emph{-\ru{ла}-}, so the same word came back from church books as \emph{\ru{глава}}, \textquotedblleft{}chapter; the head of something\textquotedblright{}.
\end{cousinweb}

\begin{grammarlens}[title={What you've built}]
One more word for this chapter.
\end{grammarlens}

\subsection*{Your turn}
\begin{itemize}
\raggedright
  \item \emph{Say it:} \emph{golová}
  \item \emph{Say it:} \emph{golová}, once more
  \item \emph{Say it:} say \emph{golová}
  \item \emph{Recall:} say the Russian for first, then the Russian for second, then say \emph{golová} again
\end{itemize}

\begin{tcolorbox}[breakable,colback=teal!4,colframe=teal!35!black,title={Before you move on}]
What does \ru{голова} mean? (\textquotedblleft{}The head\textquotedblright{}.) Say it once more.
\end{tcolorbox}
\section[\texorpdfstring{\ru{рука} (ruká)}{рука (ruká)}]{\ru{рука} (ruká) — a hand, an arm}

\label{lesson:RU-C29-hand}



\begin{quote}
Before the new one: say the Russian for second, then the Russian for the head.
\end{quote}

\subsection*{You'll want to know: \ru{рука}}
\textbf{\ru{рука}} (\emph{ruká}) — \textquotedblleft{}a hand, an arm\textquotedblright{}. One word for both.

\begin{cousinweb}
\textbf{\ru{рука}} $\leftarrow$ Proto-Slavic \emph{*rǫka}, \textquotedblleft{}hand\textquotedblright{} — inherited; the old nasal \emph{ǫ} became Russian \emph{\ru{у}}. Vasmer pairs it with Lithuanian \emph{rankà}, \textquotedblleft{}hand\textquotedblright{}, and links both with Lithuanian \emph{riñkti}, \textquotedblleft{}to gather\textquotedblright{}: the hand as the gatherer.
\end{cousinweb}

\begin{grammarlens}[title={What you've built}]
One more word for this chapter.
\end{grammarlens}

\subsection*{Your turn}
\begin{itemize}
\raggedright
  \item \emph{Say it:} \emph{ruká}
  \item \emph{Say it:} \emph{ruká}, once more
  \item \emph{Say it:} say \emph{ruká}
  \item \emph{Recall:} say the Russian for second, then the Russian for the head, then say \emph{ruká} again
\end{itemize}

\begin{tcolorbox}[breakable,colback=teal!4,colframe=teal!35!black,title={Before you move on}]
What does \ru{рука} mean? (\textquotedblleft{}A hand, an arm\textquotedblright{}.) Say it once more.
\end{tcolorbox}
\section[\texorpdfstring{\ru{нога} (nogá)}{нога (nogá)}]{\ru{нога} (nogá) — a leg, a foot}

\label{lesson:RU-C29-leg}



\begin{quote}
Before the new one: say the Russian for the head, then the Russian for a hand.
\end{quote}

\subsection*{You'll want to know: \ru{нога}}
\textbf{\ru{нога}} (\emph{nogá}) — \textquotedblleft{}a leg, a foot\textquotedblright{}. One word for both, like \textbf{\ru{рука}}.

\begin{cousinweb}
\textbf{\ru{нога}} $\leftarrow$ Proto-Slavic \emph{*noga} — inherited. Vasmer gives \textquotedblleft{}claw\textquotedblright{} as the first sense and lists Lithuanian \emph{nagà}, \textquotedblleft{}hoof\textquotedblright{}, Latin \emph{unguis}, \textquotedblleft{}nail\textquotedblright{}, and Old High German \emph{nagal}, \textquotedblleft{}nail\textquotedblright{}: English \emph{nail} is a cousin, from PIE \emph{*h\textsubscript{3}nog\textsuperscript{w}\textsuperscript{h}-}, \textquotedblleft{}nail, claw\textquotedblright{}.
\end{cousinweb}

\begin{grammarlens}[title={What you've built}]
One more word for this chapter.
\end{grammarlens}

\subsection*{Your turn}
\begin{itemize}
\raggedright
  \item \emph{Say it:} \emph{nogá}
  \item \emph{Say it:} \emph{nogá}, once more
  \item \emph{Say it:} say \emph{nogá}
  \item \emph{Recall:} say the Russian for the head, then the Russian for a hand, then say \emph{nogá} again
\end{itemize}

\begin{tcolorbox}[breakable,colback=teal!4,colframe=teal!35!black,title={Before you move on}]
What does \ru{нога} mean? (\textquotedblleft{}A leg, a foot\textquotedblright{}.) Say it once more.
\end{tcolorbox}
\section[\texorpdfstring{\ru{живот} (zhivót)}{живот (zhivót)}]{\ru{живот} (zhivót) — the belly}

\label{lesson:RU-C29-belly}



\begin{quote}
Before the new one: say the Russian for a hand, then the Russian for a leg.
\end{quote}

\subsection*{You'll want to know: \ru{живот}}
\textbf{\ru{живот}} (\emph{zhivót}) — \textquotedblleft{}the belly\textquotedblright{}.

\begin{cousinweb}
\textbf{\ru{живот}} $\leftarrow$ Proto-Slavic \emph{*životŭ}, \textquotedblleft{}life\textquotedblright{} — inherited, from the root of \emph{\ru{жить}}, \textquotedblleft{}to live\textquotedblright{}. Old Russian \emph{životŭ} meant \textquotedblleft{}life, property\textquotedblright{}; \textquotedblleft{}belly\textquotedblright{} came later (Vasmer). Lithuanian \emph{gyvatà}, \textquotedblleft{}life\textquotedblright{}, and Latin \emph{vīta}, \textquotedblleft{}life\textquotedblright{}, are cousins (Vasmer).
\end{cousinweb}

\begin{grammarlens}[title={What you've built}]
One more word for this chapter.
\end{grammarlens}

\subsection*{Your turn}
\begin{itemize}
\raggedright
  \item \emph{Say it:} \emph{zhivót}
  \item \emph{Say it:} \emph{zhivót}, once more
  \item \emph{Say it:} say \emph{zhivót}
  \item \emph{Recall:} say the Russian for a hand, then the Russian for a leg, then say \emph{zhivót} again
\end{itemize}

\begin{tcolorbox}[breakable,colback=teal!4,colframe=teal!35!black,title={Before you move on}]
What does \ru{живот} mean? (\textquotedblleft{}The belly\textquotedblright{}.) Say it once more.
\end{tcolorbox}
\section[\texorpdfstring{\ru{зуб} (zub)}{зуб (zub)}]{\ru{зуб} (zub) — a tooth}

\label{lesson:RU-C29-tooth}



\begin{quote}
Before the new one: say the Russian for a leg, then the Russian for the belly.
\end{quote}

\subsection*{You'll want to know: \ru{зуб}}
\textbf{\ru{зуб}} (\emph{zub}) — \textquotedblleft{}a tooth\textquotedblright{}.

\begin{cousinweb}
\textbf{\ru{зуб}} $\leftarrow$ Proto-Slavic \emph{*zǫbŭ}, \textquotedblleft{}tooth\textquotedblright{} — inherited. Vasmer sets it beside Sanskrit \emph{jámbha-}, \textquotedblleft{}tooth\textquotedblright{}, Greek \emph{gómphos}, \textquotedblleft{}peg\textquotedblright{}, and Old High German \emph{kamb}, \textquotedblleft{}comb\textquotedblright{}, so English \emph{comb} is a cousin, from PIE \emph{*ǵómb\textsuperscript{h}os}, \textquotedblleft{}tooth, peg\textquotedblright{}.
\end{cousinweb}

\begin{grammarlens}[title={What you've built}]
One more word for this chapter.
\end{grammarlens}

\subsection*{Your turn}
\begin{itemize}
\raggedright
  \item \emph{Say it:} \emph{zub}
  \item \emph{Say it:} \emph{zub}, once more
  \item \emph{Say it:} say \emph{zub}
  \item \emph{Recall:} say the Russian for a leg, then the Russian for the belly, then say \emph{zub} again
\end{itemize}

\begin{tcolorbox}[breakable,colback=teal!4,colframe=teal!35!black,title={Before you move on}]
What does \ru{зуб} mean? (\textquotedblleft{}A tooth\textquotedblright{}.) Say it once more.
\end{tcolorbox}
